\chapter{Varianza e disuguaglianze}

La media dice \emph{dove} si trova una variabile aleatoria, ma non \emph{quanto} è sparpagliata attorno a quel punto.
Due variabili possono avere la stessa media ed essere diversissime: una sempre vicina alla media, l'altra capace di
allontanarsene molto. In questo capitolo introduciamo le misure della dispersione, il \textbf{valore quadratico
medio}, la \textbf{varianza} e la \textbf{deviazione standard}; le calcoliamo per la binomiale, l'uniforme e la Poisson
(e per farlo ci serviranno le \textbf{derivate}, spiegate da zero); infine, con le disuguaglianze di \textbf{Markov} e
di \textbf{Chebyshev}, scopriamo che cosa ci dice davvero la deviazione standard.

% ============================================================
\section{Misurare la dispersione}
% ============================================================

Si considerino due variabili aleatorie:
\[
  X: \quad P(X = 4) = \tfrac14,\ P(X = 5) = \tfrac12,\ P(X = 6) = \tfrac14 ;
  \qquad
  Y: \quad P(Y = 1) = \tfrac14,\ P(Y = 5) = \tfrac12,\ P(Y = 9) = \tfrac14 .
\]
Entrambe hanno media 5 (le pmf sono simmetriche attorno a 5), ma $Y$ si allontana dalla media molto più di $X$.

\begin{figure}[H]
\centering
\begin{tikzpicture}
\begin{groupplot}[group style={group size=2 by 1, horizontal sep=1.4cm},
  width=6.4cm, height=4.2cm, ymin=0, ymax=0.62, xmin=0, xmax=10, xtick={1,...,9},
  ytick={0.25,0.5}, yticklabels={$\frac14$,$\frac12$}, axis lines=left, tick label style={font=\scriptsize},
  title style={font=\small}, clip=false]
\nextgroupplot[title={$X$: poco dispersa}]
  \addplot[bastoncini, mspiA] coordinates {(4,0.25) (5,0.5) (6,0.25)};
  \draw[mspiB, dashed] (axis cs:5,0) -- (axis cs:5,0.6) node[above, font=\scriptsize] {$\mu = 5$};
\nextgroupplot[title={$Y$: molto dispersa}]
  \addplot[bastoncini, mspiA] coordinates {(1,0.25) (5,0.5) (9,0.25)};
  \draw[mspiB, dashed] (axis cs:5,0) -- (axis cs:5,0.6) node[above, font=\scriptsize] {$\mu = 5$};
\end{groupplot}
\end{tikzpicture}
\caption{Stessa media, dispersione molto diversa.}
\end{figure}

L'idea naturale è misurare, in media, quanto la variabile si allontana dalla sua media $\mu_X$. Ma la media dello
scarto $X - \mu_X$ non funziona: gli scarti positivi e quelli negativi si compensano, e il risultato è sempre zero,
$E[X - \mu_X] = E[X] - \mu_X = 0$ (lo dimostreremo fra poco, con la linearità della media). Per evitare le
compensazioni si rendono positivi tutti gli scarti \textbf{elevandoli al quadrato}.

% ============================================================
\section{Valore quadratico medio e varianza}
% ============================================================

Data una variabile aleatoria $X \sim p_X(x)$, $x \in \mathcal{X}$, con media $\mu_X = E[X]$, definiamo quattro
grandezze. Tutte si calcolano con il teorema fondamentale (Teorema~\ref{teo:valoremedio}), cioè come medie di una
funzione di $X$.

\dfn{Valore quadratico medio e valore efficace}{
  Il \term{valore quadratico medio} (\emph{mean square}) di $X$ è la media del suo quadrato:
  \[ X_{rms}^2 = E[X^2] = \sum_{x \in \mathcal{X}} x^2\, p_X(x) . \]
  Il \term{valore efficace} (\emph{root mean square}) è la sua radice:
  $X_{rms} = \sqrt{E[X^2]} = \sqrt{\sum_{x \in \mathcal{X}} x^2\, p_X(x)}$.
}

\dfn{Varianza e deviazione standard}{
  La \term{varianza} di $X$ è la media del quadrato dello scarto dalla media:
  \[ \sigma_X^2 = E\left[(X - \mu_X)^2\right] = \sum_{x \in \mathcal{X}} (x - \mu_X)^2\, p_X(x) . \]
  La \term{deviazione standard} è la sua radice: $\sigma_X = \sqrt{E[(X - \mu_X)^2]} =
  \sqrt{\sum_{x} (x - \mu_X)^2\, p_X(x)}$.
}

La deviazione standard ha la stessa unità di misura di $X$ (se $X$ è in metri, $\sigma_X$ è in metri, mentre
$\sigma_X^2$ è in metri quadrati): per questo è la misura di dispersione che si usa per «leggere» i risultati.

\teorema[varianza]{Varianza, valore quadratico medio e media}{
  \[ \sigma_X^2 = X_{rms}^2 - \mu_X^2 = E[X^2] - E^2[X] , \qquad\qquad
     \sigma_X = \sqrt{X_{rms}^2 - \mu_X^2} = \sqrt{E[X^2] - E^2[X]} , \]
  dove $E^2[X]$ indica $(E[X])^2$.
}

\begin{dimostrazione}
\passo{Si sviluppa il quadrato} $(x - \mu_X)^2 = x^2 - 2 x \mu_X + \mu_X^2$, quindi
\[ \sigma_X^2 = \sum_{x \in \mathcal{X}} \left( x^2 - 2 x \mu_X + \mu_X^2 \right) p_X(x) . \]
\passo{Si spezza la somma in tre} Le costanti $2\mu_X$ e $\mu_X^2$ escono dalle sommatorie:
\[ \sigma_X^2 = \sum_{x} x^2\, p_X(x) - 2\mu_X \sum_{x} x\, p_X(x) + \mu_X^2 \sum_{x} p_X(x) . \]
\passo{Si riconoscono i pezzi} La prima somma è $E[X^2] = X_{rms}^2$, la seconda è $E[X] = \mu_X$, la terza è 1
(la pmf somma a 1):
\[ \sigma_X^2 = X_{rms}^2 - 2\mu_X^2 + \mu_X^2 = X_{rms}^2 - \mu_X^2 . \]
\end{dimostrazione}

\begin{formula}[varianza]
\[ \sigma_X^2 = E\left[(X - \mu_X)^2\right] = E[X^2] - \mu_X^2 \qquad\qquad \sigma_X = \sqrt{\sigma_X^2} \]
\end{formula}

\begin{esempio}{Le due variabili dell'inizio}
Per $X$: $\sigma_X^2 = (4-5)^2 \cdot \frac14 + (5-5)^2 \cdot \frac12 + (6-5)^2 \cdot \frac14 = \frac12$, quindi
$\sigma_X \approx 0{,}71$. Per $Y$: $\sigma_Y^2 = (1-5)^2 \cdot \frac14 + 0 + (9-5)^2 \cdot \frac14 = 8$, quindi
$\sigma_Y \approx 2{,}83$. La varianza ha colto esattamente la differenza che si vedeva nei grafici.

Con la seconda formula: $E[X^2] = 16 \cdot \frac14 + 25 \cdot \frac12 + 36 \cdot \frac14 = 25{,}5$ e
$\sigma_X^2 = 25{,}5 - 5^2 = 0{,}5$. Stesso risultato.
\end{esempio}

\info{Perché si chiama «valore efficace»}{
  Il nome viene dall'elettrotecnica. La potenza dissipata da un resistore $R$ con tensione $v$ è $v^2/R$: dipende dal
  \emph{quadrato} della tensione. Se la tensione varia nel tempo, la potenza media dipende dalla media di $v^2$, cioè
  dal valore quadratico medio. Il valore efficace è la tensione costante che dissiperebbe la stessa potenza media: i
  230 volt della presa di casa sono il valore efficace di una tensione alternata che oscilla fra circa $-325$ e
  $+325$ volt.
}

% ============================================================
\section{La varianza della binomiale}
% ============================================================

\subsection{La variabile di Bernoulli}

Cominciamo dal caso $N = 1$, $X \sim \mathcal{B}(1, p)$: la variabile vale 1 con probabilità $p$ e 0 con probabilità
$1 - p$.
\[
  E[X] = 1 \cdot p + 0 \cdot (1 - p) = p , \qquad
  E[X^2] = 1^2 \cdot p + 0^2 \cdot (1 - p) = p ,
\]
\[
  \sigma_X^2 = E[X^2] - E^2[X] = p - p^2 = p\,(1 - p) , \qquad \sigma_X = \sqrt{p\,(1 - p)} .
\]

\begin{figure}[H]
\centering
\begin{tikzpicture}
\begin{axis}[quaderno, width=8cm, height=5cm, xmin=0, xmax=1.12, ymin=0, ymax=0.32, xtick={0,0.5,1},
  ytick={0.25}, yticklabels={$\frac14$}, xlabel=$p$, ylabel={$p\,(1-p)$}]
  \addplot[mspiA, domain=0:1, samples=60] {x*(1-x)};
  \draw[mspiB, dashed] (axis cs:0.5,0) -- (axis cs:0.5,0.25) -- (axis cs:0,0.25);
  \fill[mspiB] (axis cs:0.5,0.25) circle (1.6pt);
\end{axis}
\end{tikzpicture}
\caption{La varianza di una Bernoulli in funzione di $p$. È nulla per $p = 0$ e $p = 1$ (il risultato è certo) ed è
massima, $\frac14$, per $p = \frac12$: la moneta onesta è la più imprevedibile.}
\end{figure}

\subsection{Il caso generale}

Per $X \sim \mathcal{B}(N, p)$ serve $E[X^2]$:
\[
  E[X^2] = \sum_{x \in \mathcal{X}} x^2\, p_X(x) = \sum_{k=0}^{N} k^2 \binom{N}{k} p^k q^{N-k}
  = \sum_{k=1}^{N} k^2\, \frac{N!}{k!\,(N-k)!}\, p^k q^{N-k}
  = N p \sum_{k=1}^{N} k\, \frac{(N-1)!}{(k-1)!\,(N-k)!}\, p^{k-1} q^{N-k} ,
\]
dove, come per la media, si è tolto il termine $k = 0$ (che è nullo), si è semplificato un $k$ con il fattoriale
($k^2/k! = k/(k-1)!$) e si sono portati fuori $N$ (da $N! = N\,(N-1)!$) e un $p$. A differenza della media, però,
resta un fattore $k$ davanti, e la somma non è più un binomio di Newton. Il trucco per toglierlo è usare una
\textbf{derivata}.

\begin{richiamo}{la derivata}
Si disegni il grafico di una funzione $f(x)$. La \term{derivata} $f'(x)$ dice quanto è \textbf{ripida} la curva nel
punto $x$: è la pendenza della retta tangente, cioè di quanto sale la curva per ogni passo verso destra.

Per calcolarla si prende un secondo punto vicino, $x + h$. La retta che passa per i due punti (la \emph{secante}) ha
pendenza «salita diviso spostamento», il \term{rapporto incrementale}:
\[ \frac{f(x + h) - f(x)}{h} . \]
Quando $h$ si avvicina a zero, il secondo punto scivola sul primo e la secante diventa la tangente. La derivata è il
limite del rapporto incrementale:
\[ f'(x) = \frac{\dd f}{\dd x} = \lim_{h \to 0} \frac{f(x + h) - f(x)}{h} . \]
\begin{center}
\begin{tikzpicture}
\begin{axis}[quaderno, width=8.5cm, height=6cm, xmin=-0.2, xmax=3.3, ymin=-0.3, ymax=4.6, xtick={1,2.5},
  xticklabels={$x$,$x+h$}, ytick=\empty, xlabel=$x$, ylabel=$f(x)$, grid=none]
  \addplot[mspigray, domain=0:3.05, samples=60] {0.45*x^2};
  \addplot[mspiB, domain=0.2:3.1] {0.45 + 1.575*(x-1)};
  \addplot[mspiA, domain=0.1:2.6] {0.45 + 0.9*(x-1)};
  \fill[black] (axis cs:1,0.45) circle (1.5pt); \fill[black] (axis cs:2.5,2.8125) circle (1.5pt);
  \draw[dashed, mspigray] (axis cs:1,0.45) -- (axis cs:2.5,0.45) -- (axis cs:2.5,2.8125);
  \node[font=\scriptsize, below] at (axis cs:1.75,0.45) {$h$};
  \node[font=\scriptsize, right] at (axis cs:2.5,1.15) {$f(x+h) - f(x)$};
  \node[font=\scriptsize, mspiB!85!black, anchor=south east] at (axis cs:2.2,3.3) {secante};
  \node[font=\scriptsize, mspiA!85!black, anchor=west] at (axis cs:2.62,2.15) {tangente};
\end{axis}
\end{tikzpicture}
\end{center}
\textbf{Esempio}: per $f(x) = x^2$ il rapporto incrementale è
$\frac{(x+h)^2 - x^2}{h} = \frac{2xh + h^2}{h} = 2x + h$, che per $h \to 0$ tende a $2x$. Quindi $(x^2)' = 2x$: la
parabola è piatta in $x = 0$ e sempre più ripida allontanandosi.

\textbf{Le regole} che ci servono:
\begin{itemize}
  \item una costante ha derivata nulla: $(c)' = 0$ (il suo grafico è piatto);
  \item $(x^n)' = n\,x^{n-1}$. Per il binomio di Newton $(x+h)^n = x^n + n\,x^{n-1}h + {}$(termini con $h^2$, $h^3$,
        \ldots); togliendo $x^n$ e dividendo per $h$ resta $n\,x^{n-1} + {}$(termini con $h$), che tendono a zero;
  \item \textbf{linearità}: $(a f + b g)' = a f' + b g'$; in particolare la derivata di una somma (anche di molti
        termini) è la somma delle derivate;
  \item \textbf{regola del prodotto}: $(f g)' = f' g + f g'$ (spiegata dal disegno qui sotto);
  \item $\left((x + c)^m\right)' = m\,(x + c)^{m-1}$: spostare il grafico in orizzontale non cambia le pendenze.
\end{itemize}
\begin{center}
\begin{tikzpicture}[x=1cm, y=1cm]
  \fill[mspiA!15] (0,0) rectangle (3,2);
  \fill[mspiB!25] (3,0) rectangle (3.7,2);
  \fill[mspiB!25] (0,2) rectangle (3,2.5);
  \fill[mspiC!35] (3,2) rectangle (3.7,2.5);
  \draw (0,0) rectangle (3.7,2.5); \draw (3,0) -- (3,2.5); \draw (0,2) -- (3.7,2);
  \node[font=\small] at (1.5,1) {$f \cdot g$};
  \node[font=\scriptsize] at (3.35,1) {$g\,\Delta f$};
  \node[font=\scriptsize] at (1.5,2.25) {$f\,\Delta g$};
  \draw[decorate, decoration={brace, amplitude=4pt, mirror}] (0,-0.1) -- (3,-0.1) node[midway, below=4pt,
    font=\scriptsize] {$f$};
  \draw[decorate, decoration={brace, amplitude=4pt, mirror}] (3,-0.1) -- (3.7,-0.1) node[midway, below=4pt,
    font=\scriptsize] {$\Delta f$};
  \draw[decorate, decoration={brace, amplitude=4pt}] (-0.1,0) -- (-0.1,2) node[midway, left=4pt,
    font=\scriptsize] {$g$};
  \draw[decorate, decoration={brace, amplitude=4pt}] (-0.1,2) -- (-0.1,2.5) node[midway, left=4pt,
    font=\scriptsize] {$\Delta g$};
  \node[font=\small, anchor=west, text width=7.6cm] at (4.3,1.25)
    {Il prodotto $f \cdot g$ è l'area di un rettangolo. Quando $x$ cresce di $h$, i lati crescono di $\Delta f$ e
     $\Delta g$, e l'area cresce di $f\,\Delta g + g\,\Delta f + \Delta f\,\Delta g$ (le strisce arancioni e il
     quadratino verde). Dividendo per $h$ e mandando $h$ a zero, $\Delta f/h \to f'$, $\Delta g/h \to g'$, mentre il
     quadratino, che è piccolo «due volte», sparisce: resta $f' g + f g'$.};
\end{tikzpicture}
\end{center}
\textbf{Derivare rispetto a una lettera tenendo ferme le altre.} Se un'espressione contiene due lettere, per esempio
$p$ e $q$, si può derivare rispetto a $p$ trattando $q$ come un numero fisso: $\frac{\dd}{\dd p}\left(p^2 q^3\right)
= 2p\, q^3$. Lo useremo con $q$, che solo alla fine sostituiremo con $1 - p$.
\end{richiamo}

\teorema[varbinomiale]{Varianza della binomiale}{
  Se $X \sim \mathcal{B}(N, p)$, allora
  \[ E[X^2] = N p\,(1 - p) + N^2 p^2 , \qquad \sigma_X^2 = N p\,(1 - p) , \qquad \sigma_X = \sqrt{N p\,(1 - p)} . \]
}

\begin{dimostrazione}[il trucco della derivata]
\passo{L'idea} Un fattore $k$ davanti a $p^{k-1}$ è proprio quello che produce la derivata: $k\,p^{k-1} =
\frac{\dd}{\dd p}\, p^k$. Trattando $q$ come una lettera indipendente da $p$ (e sostituendo $q = 1 - p$ solo alla fine),
ogni termine della somma è la derivata rispetto a $p$ di un termine senza il fattore $k$:
\[
  E[X^2] = N p \sum_{k=1}^{N} \frac{(N-1)!}{(k-1)!\,(N-k)!}\, \frac{\dd}{\dd p}\left(p^k\right) q^{N-k}
  = N p\, \frac{\dd}{\dd p} \left[ \sum_{k=1}^{N} \frac{(N-1)!}{(k-1)!\,[(N-1)-(k-1)]!}\, p^k q^{N-k} \right]_{q = 1-p}
\]
(la derivata di una somma è la somma delle derivate, e i fattori che non contengono $p$ restano fuori dalla derivata).
\passo{Si riconosce un binomio di Newton} Posto $l = k - 1$, e staccando un fattore $p$ da $p^k = p \cdot p^l$:
\[
  \sum_{k=1}^{N} \frac{(N-1)!}{(k-1)!\,[(N-1)-(k-1)]!}\, p^k q^{N-k}
  = p \sum_{l=0}^{N-1} \frac{(N-1)!}{l!\,[(N-1) - l]!}\, p^l q^{N-1-l} = p\,(p + q)^{N-1} .
\]
\passo{Si deriva} Con la regola del prodotto, e ricordando che $q$ è fisso:
\[ \frac{\dd}{\dd p}\left[ p\,(p + q)^{N-1} \right] = (p + q)^{N-1} + p\,(N-1)\,(p + q)^{N-2} . \]
\passo{Si sostituisce $q = 1 - p$} Ora $p + q = 1$, e le sue potenze valgono 1:
\[ E[X^2] = N p \left[ 1 + p\,(N-1) \right] = N p \left[ 1 + N p - p \right] = N p\,(1 - p) + N^2 p^2 . \]
\passo{La varianza} Siccome $\mu_X = N p$ (Teorema~\ref{teo:mediabinomiale}), il secondo addendo è proprio $\mu_X^2$:
\[ \sigma_X^2 = E[X^2] - \mu_X^2 = N p\,(1 - p) + N^2 p^2 - N^2 p^2 = N p\,(1 - p) . \]
\end{dimostrazione}

Lo stesso risultato si ottiene in un modo più semplice, che evita le derivate: invece di $E[X^2]$ si calcola
$E[X(X-1)]$, perché il fattore $k(k-1)$ si semplifica con il fattoriale ancora meglio di $k^2$.

\begin{dimostrazione}[secondo modo: il momento fattoriale]
\passo{Si calcola $E[X(X-1)]$} Per il teorema fondamentale, con $g(x) = x(x-1)$. I termini con $k = 0$ e $k = 1$ sono
nulli; per $k \geq 2$, $\frac{k(k-1)}{k!} = \frac{1}{(k-2)!}$:
\[
  E[X(X-1)] = \sum_{k=2}^{N} k(k-1) \frac{N!}{k!\,(N-k)!}\, p^k (1-p)^{N-k}
  = N(N-1)\,p^2 \sum_{k=2}^{N} \frac{(N-2)!}{(k-2)!\,[N-2-(k-2)]!}\, p^{k-2} (1-p)^{N-2-(k-2)} .
\]
\passo{Newton ancora una volta} Con $l = k - 2$ e $M = N - 2$ la somma diventa
$\sum_{l=0}^{M} \frac{M!}{l!\,(M-l)!}\, p^l (1-p)^{M-l} = 1$, quindi $E[X(X-1)] = N(N-1)\,p^2$.
\passo{Si torna a $E[X^2]$} Sempre per il teorema fondamentale, la media di una differenza si spezza:
$E[X(X-1)] = \sum_x (x^2 - x)\, p_X(x) = E[X^2] - E[X]$. Quindi
\[ E[X^2] = E[X(X-1)] + E[X] = N(N-1)\,p^2 + N p = N^2 p^2 - N p^2 + N p = N p\,(1 - p) + N^2 p^2 . \]
\end{dimostrazione}

\begin{formula}[binomiale $\mathcal{B}(N, p)$]
\[ \mu_X = N p \qquad X_{rms}^2 = N p\,(1-p) + N^2 p^2 \qquad \sigma_X^2 = N p\,(1-p) \qquad
   \sigma_X = \sqrt{N p\,(1-p)} \]
\end{formula}

% ============================================================
\section{La varianza dell'uniforme}
% ============================================================

Sia $X \sim \mathcal{U}(\{0, 1, \ldots, M-1\})$, per cui $\mu_X = \frac{M-1}{2}$ (Capitolo~5). Il valore quadratico
medio è
\[ X_{rms}^2 = \frac{1}{M} \sum_{k=0}^{M-1} k^2 = \frac{1}{M} \sum_{k=1}^{M-1} k^2 , \]
e serve la somma dei primi quadrati. Invece di cercarla per tentativi, la ricaviamo con un trucco che spiega da dove
viene la formula.

\begin{richiamo}{le somme telescopiche}
Una somma in cui ogni termine è una differenza fra due valori consecutivi, $a_{l+1} - a_l$, si semplifica quasi tutta:
\[
  \sum_{l=1}^{n} (a_{l+1} - a_l) = (a_2 - a_1) + (a_3 - a_2) + (a_4 - a_3) + \cdots + (a_{n+1} - a_n) = a_{n+1} - a_1 .
\]
Ogni $a_l$ intermedio compare una volta con il più e una volta con il meno, e si cancella: restano solo l'ultimo e il
primo. Il nome viene dal cannocchiale (telescopio) che si chiude, con i tubi che rientrano uno nell'altro.
\end{richiamo}

\teorema[quadrati]{Somma dei primi quadrati}{
  Per ogni naturale $n \geq 1$:
  \[ \sum_{l=1}^{n} l^2 = 1^2 + 2^2 + \cdots + n^2 = \frac{n\,(n+1)\,(2n+1)}{6} . \]
}

\begin{dimostrazione}[con una somma telescopica]
\passo{Il cubo di un binomio} Per ogni $l$: $(l+1)^3 = l^3 + 3l^2 + 3l + 1$, quindi
\[ (l+1)^3 - l^3 = 3l^2 + 3l + 1 . \]
\passo{Si somma da 1 a $n$} A sinistra c'è una somma telescopica, che vale $(n+1)^3 - 1^3$. A destra si spezza la
somma in tre, chiamando $S$ la somma dei quadrati che cerchiamo e usando la somma di Gauss
(Teorema~\ref{teo:gauss}):
\[ (n+1)^3 - 1 = 3S + 3\, \frac{n\,(n+1)}{2} + n . \]
\passo{Si ricava $S$} Si porta tutto il resto a sinistra e si raccoglie $(n+1)$:
\[
  3S = (n+1)^3 - 1 - n - \frac{3n\,(n+1)}{2} = (n+1)^3 - (n+1) - \frac{3n\,(n+1)}{2}
     = (n+1) \left[ (n+1)^2 - 1 - \frac{3n}{2} \right] .
\]
Dentro la parentesi: $(n+1)^2 - 1 - \frac{3n}{2} = n^2 + 2n - \frac{3n}{2} = n^2 + \frac{n}{2} = \frac{n\,(2n+1)}{2}$.
Quindi $3S = \frac{n\,(n+1)\,(2n+1)}{2}$ e $S = \frac{n\,(n+1)\,(2n+1)}{6}$.
\passo{Una verifica} Con $n = 3$: $1 + 4 + 9 = 14$ e $\frac{3 \cdot 4 \cdot 7}{6} = 14$.
\end{dimostrazione}

\teorema[varuniforme]{Varianza dell'uniforme}{
  Se $X \sim \mathcal{U}(\{0, 1, \ldots, M-1\})$:
  \[ \mu_X = \frac{M-1}{2} , \qquad X_{rms}^2 = \frac{(M-1)(2M-1)}{6} , \qquad \sigma_X^2 = \frac{M^2 - 1}{12} . \]
}

\begin{dimostrazione}
\passo{Il valore quadratico medio} Con la somma dei quadrati, $n = M - 1$:
\[ X_{rms}^2 = \frac{1}{M} \cdot \frac{(M-1)\,M\,(2M-1)}{6} = \frac{(M-1)(2M-1)}{6} . \]
\passo{La varianza} Si sottrae $\mu_X^2 = (M-1)^2/4$ e si porta tutto al denominatore comune 12:
\[
  \sigma_X^2 = \frac{(M-1)(2M-1)}{6} - \frac{(M-1)^2}{4} = \frac{2(M-1)(2M-1) - 3(M-1)^2}{12} .
\]
\passo{Si sviluppano i prodotti}
\[
  \sigma_X^2 = \frac{2\,(2M^2 - M - 2M + 1) - 3\,(M^2 - 2M + 1)}{12}
  = \frac{4M^2 - 6M + 2 - 3M^2 + 6M - 3}{12} = \frac{M^2 - 1}{12} .
\]
\end{dimostrazione}

\begin{esempio}{Dado e byte}
\begin{itemize}
  \item Un byte casuale ($M = 256$): $\sigma_X^2 = \frac{65\,536 - 1}{12} \approx 5\,461$ e $\sigma_X \approx 73{,}9$.
  \item Un dado onesto ha valori $1, \ldots, 6$: sono i valori $0, \ldots, 5$ ($M = 6$) aumentati di 1. Vedremo fra
        poco che aggiungere una costante non cambia la varianza, quindi $\sigma_X^2 = \frac{36 - 1}{12} = \frac{35}{12}
        \approx 2{,}92$ e $\sigma_X \approx 1{,}71$.
\end{itemize}
\end{esempio}

% ============================================================
\section{La varianza della Poisson}
% ============================================================

Per la Poisson useremo lo stesso trucco della derivata, questa volta su una serie. Serve sapere come si deriva
l'esponenziale; è anche l'occasione per capire, finalmente, da dove viene la serie esponenziale.

\begin{richiamo}{la derivata dell'esponenziale e la serie esponenziale}
\textbf{La derivata di $e^x$ è $e^x$.} Il rapporto incrementale è
\[ \frac{e^{x+h} - e^x}{h} = e^x\, \frac{e^h - 1}{h} , \]
e il secondo fattore tende a 1: è la proprietà che distingue il numero $e$ da tutte le altre basi.
\begin{center}
\renewcommand{\arraystretch}{1.2}
\begin{tabular}{@{}c|cccc@{}}
  $h$ & 0,1 & 0,01 & 0,001 & 0,0001 \\ \hline
  $\dfrac{e^h - 1}{h}$ & 1,0517 & 1,0050 & 1,0005 & 1,00005
\end{tabular}
\end{center}
Quindi $(e^x)' = e^x$: la curva $e^x$ è ripida esattamente quanto è alta (con la base 2 la pendenza in 0 sarebbe circa
$0{,}69$, con la base 3 circa $1{,}10$; con $e$ è esattamente 1).

\textbf{La derivata di $e^{-x}$.} Siccome $e^{-x} \cdot e^{x} = 1$, derivando con la regola del prodotto si ottiene
$(e^{-x})'\,e^x + e^{-x}\,e^x = 0$, cioè $(e^{-x})' = -e^{-x}$.

\textbf{Da dove viene la serie.} Sia $S(x) = \sum_{k=0}^{\infty} \frac{x^k}{k!} = 1 + x + \frac{x^2}{2} + \cdots$.
Una serie di potenze come questa si può derivare termine a termine, e
$\left(\frac{x^k}{k!}\right)' = \frac{k\,x^{k-1}}{k!} = \frac{x^{k-1}}{(k-1)!}$: ogni termine diventa il precedente,
quindi $S'(x) = S(x)$. Inoltre $S(0) = 1$. Anche $e^x$ ha queste due proprietà, e le due funzioni coincidono: la
funzione $D(x) = S(x)\, e^{-x}$ ha derivata
\[ D'(x) = S'(x)\, e^{-x} + S(x)\, (e^{-x})' = S(x)\, e^{-x} - S(x)\, e^{-x} = 0 , \]
quindi è costante (una curva con pendenza sempre nulla è piatta) e vale $D(0) = 1$. Allora $S(x)\,e^{-x} = 1$, cioè
$S(x) = e^x$.
\end{richiamo}

\teorema[varpoisson]{Varianza della Poisson}{
  Se $X \sim \mathcal{P}(\lambda)$:
  \[ \mu_X = \lambda , \qquad X_{rms}^2 = \lambda + \lambda^2 , \qquad \sigma_X^2 = \lambda , \qquad
     \sigma_X = \sqrt{\lambda} . \]
  Per una variabile di Poisson la media e la varianza coincidono.
}

\begin{dimostrazione}[il trucco della derivata]
\passo{Si prepara la somma} Si toglie il termine $k = 0$ (nullo), si semplifica $k^2/k! = k/(k-1)!$ e si porta fuori
$e^{-\lambda}$ e un $\lambda$:
\[
  X_{rms}^2 = e^{-\lambda} \sum_{k=0}^{\infty} k^2\, \frac{\lambda^k}{k!} = e^{-\lambda} \sum_{k=1}^{\infty}
  k\, \frac{\lambda^k}{(k-1)!} = e^{-\lambda}\, \lambda \sum_{k=1}^{\infty} k\, \frac{\lambda^{k-1}}{(k-1)!} .
\]
\passo{Il fattore $k$ diventa una derivata} $k\,\lambda^{k-1} = \frac{\dd}{\dd\lambda} \lambda^k$, e la serie si deriva
termine a termine:
\[ X_{rms}^2 = e^{-\lambda}\, \lambda\, \frac{\dd}{\dd\lambda} \left( \sum_{k=1}^{\infty} \frac{\lambda^k}{(k-1)!} \right)
   = e^{-\lambda}\, \lambda\, \frac{\dd}{\dd\lambda} \left( \lambda \sum_{k=1}^{\infty} \frac{\lambda^{k-1}}{(k-1)!}
   \right) . \]
\passo{Si riconosce l'esponenziale} Posto $l = k - 1$, la serie è $\sum_{l=0}^{\infty} \frac{\lambda^l}{l!} =
e^{\lambda}$:
\[ X_{rms}^2 = e^{-\lambda}\, \lambda\, \frac{\dd}{\dd\lambda} \left( \lambda\, e^{\lambda} \right) . \]
\passo{Si deriva con la regola del prodotto} $\frac{\dd}{\dd\lambda}\left(\lambda\, e^{\lambda}\right) = e^{\lambda} +
\lambda\, e^{\lambda}$, quindi
\[ X_{rms}^2 = e^{-\lambda}\, \lambda \left( e^{\lambda} + \lambda\, e^{\lambda} \right) = \lambda + \lambda^2 . \]
\passo{La varianza} $\sigma_X^2 = X_{rms}^2 - \mu_X^2 = \lambda + \lambda^2 - \lambda^2 = \lambda$.
\end{dimostrazione}

\begin{dimostrazione}[secondo modo: il momento fattoriale]
\passo{} I termini con $k = 0$ e $k = 1$ sono nulli; per $k \geq 2$ si semplifica $k(k-1)/k! = 1/(k-2)!$ e, posto
$l = k - 2$:
\[
  E[X(X-1)] = \sum_{k=2}^{\infty} k\,(k-1)\, e^{-\lambda} \frac{\lambda^k}{k!}
  = e^{-\lambda}\, \lambda^2 \sum_{k=2}^{\infty} \frac{\lambda^{k-2}}{(k-2)!}
  = e^{-\lambda}\, \lambda^2 \sum_{l=0}^{\infty} \frac{\lambda^l}{l!} = e^{-\lambda}\, \lambda^2\, e^{\lambda} = \lambda^2 .
\]
\passo{} Quindi $E[X^2] = E[X(X-1)] + E[X] = \lambda^2 + \lambda$.
\end{dimostrazione}

% ============================================================
\section{Le proprietà della media e della varianza}
% ============================================================

\teorema[lineare]{Media e varianza di una trasformazione lineare}{
  Siano $a$ e $b$ costanti reali e $Y = aX + b$. Allora:
  \begin{enumerate}[label=\alph*)]
    \item $E[aX + b] = a\,E[X] + b$ (in particolare $E[b] = b$: la media di una costante è la costante);
    \item $\sigma_Y^2 = a^2 \sigma_X^2$: la costante $b$ non modifica la varianza, e la costante $a$ va elevata al
          quadrato (quindi $\sigma_Y = \abs{a}\,\sigma_X$);
    \item $E[Y^2] = Y_{rms}^2 = a^2 X_{rms}^2 + b^2 + 2ab\,\mu_X$.
  \end{enumerate}
}

\begin{dimostrazione}
\passo{La media} Per il teorema fondamentale, con $g(x) = ax + b$:
\[ E[aX + b] = \sum_{x} (ax + b)\, p_X(x) = a \sum_{x} x\, p_X(x) + b \sum_{x} p_X(x) = a\,E[X] + b . \]
\passo{La varianza} Per il passo precedente $\mu_Y = a\mu_X + b$, quindi lo scarto di $Y$ dalla sua media è
$Y - \mu_Y = aX + b - a\mu_X - b = a\,(X - \mu_X)$: la $b$ si cancella. Allora
\[ \sigma_Y^2 = E\left[(Y - \mu_Y)^2\right] = E\left[a^2 (X - \mu_X)^2\right] = a^2\, E\left[(X - \mu_X)^2\right] =
   a^2 \sigma_X^2 , \]
dove si è usato il passo 1 per portare fuori la costante $a^2$.
\passo{Il valore quadratico medio} Per il Teorema~\ref{teo:varianza}, $E[Y^2] = \sigma_Y^2 + \mu_Y^2$:
\[
  E[Y^2] = a^2 \sigma_X^2 + (a\mu_X + b)^2 = a^2 \sigma_X^2 + a^2 \mu_X^2 + b^2 + 2ab\,\mu_X
  = a^2 \left(\sigma_X^2 + \mu_X^2\right) + b^2 + 2ab\,\mu_X = a^2 X_{rms}^2 + b^2 + 2ab\,\mu_X .
\]
\end{dimostrazione}

\begin{figure}[H]
\centering
\begin{tikzpicture}
\begin{groupplot}[group style={group size=3 by 1, horizontal sep=0.8cm},
  width=5.4cm, height=3.9cm, ymin=0, ymax=0.62, xmin=-0.6, xmax=6.6, xtick={0,...,6},
  ytick={0.25,0.5}, yticklabels={$\frac14$,$\frac12$}, axis lines=left, tick label style={font=\scriptsize},
  title style={font=\small}]
\nextgroupplot[title={$X$: $\mu = 1$, $\sigma^2 = \frac12$}]
  \addplot[bastoncini, mspiA] coordinates {(0,0.25) (1,0.5) (2,0.25)};
\nextgroupplot[title={$X + 3$: $\mu = 4$, $\sigma^2 = \frac12$}]
  \addplot[bastoncini, mspiB] coordinates {(3,0.25) (4,0.5) (5,0.25)};
\nextgroupplot[title={$2X$: $\mu = 2$, $\sigma^2 = 2$}]
  \addplot[bastoncini, mspiC] coordinates {(0,0.25) (2,0.5) (4,0.25)};
\end{groupplot}
\end{tikzpicture}
\caption{Aggiungere una costante sposta la pmf ma non la allarga: la varianza resta la stessa. Moltiplicare per 2
allarga le distanze di 2 volte, quindi la varianza (che è una media di distanze al quadrato) di 4 volte.}
\end{figure}

Altre tre proprietà discendono direttamente dalle definizioni:
\begin{itemize}
  \item se $X(\omega) \geq 0$ per ogni $\omega \in \Omega$ (una variabile aleatoria \textbf{non negativa}), allora
        $E[X] = \sum_x x\, p_X(x) \geq 0$, perché è una somma di prodotti di quantità non negative;
  \item $\sigma_X^2 \geq 0$ sempre, perché $\sigma_X^2 = \sum_x (x - \mu_X)^2 p_X(x)$ è una somma di prodotti di
        quantità non negative (o, se si preferisce, la media della variabile non negativa $(X - \mu_X)^2$);
  \item per le variabili a \textbf{media nulla} ($\mu_X = 0$) varianza e valore quadratico medio coincidono:
        $\sigma_X^2 = X_{rms}^2$.
\end{itemize}
In particolare, ponendo $a = 1$ e $b = -\mu_X$ nella proprietà a): $E[X - \mu_X] = \mu_X - \mu_X = 0$, come avevamo
anticipato all'inizio del capitolo.

\begin{table}[H]
\centering
\renewcommand{\arraystretch}{1.6}
\begin{tabular}{@{}lllll@{}}
  \toprule
  \textbf{Variabile} & \textbf{Alfabeto} & \textbf{pmf} & \textbf{Media} & \textbf{Varianza} \\
  \midrule
  Bernoulli $\mathcal{B}(1, p)$ & $\{0, 1\}$ & $p_X(1) = p$, $p_X(0) = 1 - p$ & $p$ & $p\,(1-p)$ \\
  binomiale $\mathcal{B}(N, p)$ & $\{0, \ldots, N\}$ & $\binom{N}{k} p^k (1-p)^{N-k}$ & $N p$ & $N p\,(1-p)$ \\
  uniforme $\mathcal{U}(\{0, \ldots, M-1\})$ & $\{0, \ldots, M-1\}$ & $\frac{1}{M}$ & $\frac{M-1}{2}$ &
    $\frac{M^2-1}{12}$ \\
  Poisson $\mathcal{P}(\lambda)$ & $\N_0$ & $\frac{\lambda^k}{k!} e^{-\lambda}$ & $\lambda$ & $\lambda$ \\
  \bottomrule
\end{tabular}
\caption{Le variabili aleatorie discrete notevoli.}
\end{table}

% ============================================================
\section{La disuguaglianza di Markov}
% ============================================================

Media e varianza sono «riassunti» della pmf. Quanto si può dire su una variabile aleatoria conoscendo solo quelli?
Sorprendentemente molto, grazie a due disuguaglianze che valgono per \textbf{qualsiasi} pmf.

\teorema[markov]{Disuguaglianza di Markov}{
  Sia $Z$ una variabile aleatoria \textbf{non negativa}, con alfabeto discreto $\mathcal{Z} \subseteq [0, +\infty[$ e
  pmf $p_Z(z)$. Allora, per qualunque $\delta > 0$:
  \[ P(Z \geq \delta) \leq \frac{E[Z]}{\delta} . \]
}

\begin{dimostrazione}
\passo{La probabilità della coda} L'evento $\{Z \geq \delta\}$ si verifica quando $Z$ assume uno dei valori
$z \geq \delta$:
\[ P(Z \geq \delta) = \sum_{z \,:\, z \geq \delta} p_Z(z) . \]
\passo{(a) Si moltiplica per un numero $\geq 1$} Per i valori $z \geq \delta$ della somma vale $\frac{z}{\delta} \geq
1$; moltiplicando ogni termine per $\frac{z}{\delta}$, la somma non può diminuire:
\[ \sum_{z \,:\, z \geq \delta} p_Z(z) \leq \sum_{z \,:\, z \geq \delta} \frac{z}{\delta}\, p_Z(z) . \]
\passo{(b) Si estende la somma a tutto l'alfabeto} Aggiungendo anche i valori $z < \delta$ si aggiungono termini
$\frac{z}{\delta} p_Z(z) \geq 0$ (perché $Z$ è non negativa): la somma non può diminuire.
\[ \sum_{z \,:\, z \geq \delta} \frac{z}{\delta}\, p_Z(z) \leq \sum_{z \in \mathcal{Z}} \frac{z}{\delta}\, p_Z(z) . \]
\passo{(c) Si riconosce la media} $\sum_{z \in \mathcal{Z}} \frac{z}{\delta}\, p_Z(z) = \frac{1}{\delta}
\sum_{z} z\, p_Z(z) = \frac{E[Z]}{\delta}$, per la definizione di media. Mettendo in fila i tre passaggi:
\[ P(Z \geq \delta) \underset{(a)}{\leq} \sum_{z \geq \delta} \frac{z}{\delta}\, p_Z(z) \underset{(b)}{\leq}
   \sum_{z \in \mathcal{Z}} \frac{z}{\delta}\, p_Z(z) \underset{(c)}{=} \frac{E[Z]}{\delta} . \]
\end{dimostrazione}

\begin{figure}[H]
\centering
\begin{tikzpicture}
\begin{axis}[width=11cm, height=4.8cm, xmin=-0.6, xmax=15.6, ymin=0, ymax=0.23, xtick={0,...,15}, axis lines=left,
  tick label style={font=\scriptsize}, xlabel=$z$, ylabel={$p_Z(z)$}, clip=false,
  yticklabel style={/pgf/number format/fixed, /pgf/number format/use comma},
  grid=major, grid style={mspigray!15}, xmajorgrids=false]
  \addplot[bastoncini, mspiA] coordinates {(0,0.0183) (1,0.0733) (2,0.1465) (3,0.1954) (4,0.1954) (5,0.1563) (6,0.1042) (7,0.0595)};
  \addplot[bastoncini, mspiB] coordinates {(8,0.0298) (9,0.0132) (10,0.0053) (11,0.0019) (12,0.0006) (13,0.0002) (14,0.0001) (15,0.0000)};
  \draw[mspigray, dashed] (axis cs:7.5,0) -- (axis cs:7.5,0.21) node[above, font=\scriptsize] {$\delta = 8$};
  \draw[mspiB, dashed] (axis cs:4,0) -- (axis cs:4,0.21) node[above, font=\scriptsize] {$E[Z] = 4$};
  \node[font=\small, anchor=west, text width=5cm, align=left] at (axis cs:9.2,0.14)
    {coda $\{Z \geq 8\}$, in arancione:\\ vera probabilità $\approx 0{,}051$\\ Markov: $\leq 4/8 = 0{,}5$};
\end{axis}
\end{tikzpicture}
\caption{Una Poisson con $\lambda = 4$. Markov garantisce che la coda oltre $\delta = 8$ abbia probabilità al più
$0{,}5$; quella vera è molto più piccola. La disuguaglianza è grossolana, ma vale per qualsiasi variabile non negativa
di cui si conosca solo la media.}
\end{figure}

\begin{esempio}{La coda del router}
Si sa soltanto che il numero medio di pacchetti in coda in un router è 2. Qual è, al massimo, la probabilità di
trovarne 10 o più? Il numero di pacchetti è una variabile non negativa, quindi per Markov
$P(Z \geq 10) \leq \frac{2}{10} = 0{,}2$. Senza sapere nient'altro della pmf, non si può dire di più; se poi si sapesse
che $Z$ è di Poisson, il conto esatto darebbe circa $0{,}00005$.
\end{esempio}

% ============================================================
\section{La disuguaglianza di Chebyshev}
% ============================================================

La disuguaglianza di Markov vale per variabili non negative. Applicandola alla variabile non negativa
$(X - \mu_X)^2$, che misura la distanza al quadrato dalla media, si ottiene un risultato che vale per
\textbf{qualsiasi} variabile e che coinvolge la deviazione standard.

\teorema[chebyshev]{Disuguaglianza di Chebyshev}{
  Sia $X$ una variabile aleatoria con media $\mu_X$ e deviazione standard $\sigma_X > 0$. Per ogni $k > 0$:
  \[ P\left( \abs{X - \mu_X} \geq k\,\sigma_X \right) \leq \frac{1}{k^2} . \]
}

\begin{dimostrazione}
\passo{Le scelte} Si applica Markov con $Z = \abs{X - \mu_X}^2$, che è non negativa, e $\delta = (k\,\sigma_X)^2 > 0$.
\passo{La media di $Z$} Per variabili reali $\abs{X - \mu_X}^2 = (X - \mu_X)^2$, quindi
\[ E[Z] = E\left[\abs{X - \mu_X}^2\right] = E\left[(X - \mu_X)^2\right] = \sigma_X^2 . \]
\passo{L'evento} Due numeri non negativi stanno nello stesso ordine dei loro quadrati, quindi
$\abs{X - \mu_X}^2 \geq (k\,\sigma_X)^2$ esattamente quando $\abs{X - \mu_X} \geq k\,\sigma_X$.
\passo{Markov} Mettendo insieme:
\[ P\left( \abs{X - \mu_X}^2 \geq (k\,\sigma_X)^2 \right) = P\left( \abs{X - \mu_X} \geq k\,\sigma_X \right) \leq
   \frac{\sigma_X^2}{k^2 \sigma_X^2} = \frac{1}{k^2} . \]
\end{dimostrazione}

% ============================================================
\section{Il significato della varianza e della deviazione standard}
% ============================================================

La disuguaglianza di Chebyshev si può leggere «al contrario», guardando la probabilità di restare \emph{vicini} alla
media. Serve un richiamo sul valore assoluto.

\begin{richiamo}{valore assoluto e intervalli}
Il \term{valore assoluto} $\abs{a}$ è il numero $a$ senza segno: $\abs{3} = \abs{-3} = 3$. La quantità
$\abs{x - \mu}$ è la \textbf{distanza} fra $x$ e $\mu$ sulla retta dei numeri. Dire che la distanza non supera $r$ vuol
dire che $x$ sta nell'intervallo centrato in $\mu$ di raggio $r$:
\[ \abs{x - \mu} \leq r \iff -r \leq x - \mu \leq r \iff \mu - r \leq x \leq \mu + r . \]
\begin{center}
\begin{tikzpicture}
  \draw[-{Stealth}] (-0.5,0) -- (8.5,0);
  \fill[evid] (2,-0.12) rectangle (6,0.12);
  \draw[line width=1pt] (2,-0.15) -- (2,0.15); \draw[line width=1pt] (6,-0.15) -- (6,0.15);
  \fill (4,0) circle (1.6pt);
  \node[font=\small, below=4pt] at (4,0) {$\mu$};
  \node[font=\small, below=4pt] at (2,0) {$\mu - r$};
  \node[font=\small, below=4pt] at (6,0) {$\mu + r$};
  \draw[decorate, decoration={brace, amplitude=4pt}] (4,0.25) -- (6,0.25) node[midway, above=4pt, font=\scriptsize]
    {distanza $r$};
  \node[font=\scriptsize, anchor=west] at (6.3,0.45) {qui $\abs{x - \mu} > r$};
\end{tikzpicture}
\end{center}
\end{richiamo}

\teorema[sigmasignificato]{Quanto si resta vicini alla media}{
  Per ogni variabile aleatoria e ogni $k > 0$:
  \[ P\left( \mu_X - k\,\sigma_X \leq X \leq \mu_X + k\,\sigma_X \right) \geq 1 - \frac{1}{k^2} . \]
}

\begin{dimostrazione}
\passo{Si passa al complementare} L'evento «distanza minore di $k\sigma_X$» è il complementare di «distanza almeno
$k\sigma_X$»:
\[ P\left(\abs{X - \mu_X} < k\,\sigma_X\right) = 1 - P\left(\abs{X - \mu_X} \geq k\,\sigma_X\right) \geq 1 -
   \frac{1}{k^2} , \]
perché togliendo da 1 un numero che è al più $1/k^2$ si ottiene almeno $1 - 1/k^2$.
\passo{Si aggiunge l'uguaglianza} Se nella disuguaglianza a sinistra si ammette anche l'uguaglianza
($\abs{X - \mu_X} \leq k\,\sigma_X$) si aggiunge una quantità non negativa (che può anche essere nulla), e la
disuguaglianza resta vera a maggior ragione.
\passo{Si toglie il valore assoluto} Per il richiamo, $\abs{X - \mu_X} \leq k\,\sigma_X$ equivale a
$-k\sigma_X \leq X - \mu_X \leq k\sigma_X$, cioè a $\mu_X - k\sigma_X \leq X \leq \mu_X + k\sigma_X$.
\end{dimostrazione}

Il teorema dice che la deviazione standard è il \textbf{righello naturale} della variabile: qualunque sia la pmf,
almeno il 75\% della probabilità sta entro 2 deviazioni standard dalla media, almeno l'89\% entro 3, almeno il 99\%
entro 10.
\begin{center}
\renewcommand{\arraystretch}{1.25}
\begin{tabular}{@{}c|cccccc@{}}
  $k$ & 1 & 1,5 & 2 & 3 & 4 & 10 \\ \hline
  $1 - 1/k^2$ & 0 & 0,56 & 0,75 & 0,89 & 0,94 & 0,99
\end{tabular}
\end{center}

\begin{figure}[H]
\centering
\begin{tikzpicture}
\begin{axis}[width=12cm, height=5cm, xmin=-0.6, xmax=20.6, ymin=0, ymax=0.23, xtick={0,2,...,20}, axis lines=left,
  tick label style={font=\scriptsize}, xlabel=$k$, ylabel={$p_X(k)$}, clip=false,
  yticklabel style={/pgf/number format/fixed, /pgf/number format/use comma},
  grid=major, grid style={mspigray!15}, xmajorgrids=false]
  \fill[evid, opacity=0.55] (axis cs:1.901,0) rectangle (axis cs:10.099,0.2);
  \addplot[bastoncini, mspiA] coordinates {(0,0.0008) (1,0.0068) (2,0.0278) (3,0.0716) (4,0.1304) (5,0.1789) (6,0.1916) (7,0.1643) (8,0.1144) (9,0.0654) (10,0.0308) (11,0.0120) (12,0.0039) (13,0.0010) (14,0.0002) (15,0.0000) (16,0.0000) (17,0.0000) (18,0.0000) (19,0.0000) (20,0.0000)};
  \draw[mspiB, dashed] (axis cs:6,0) -- (axis cs:6,0.215) node[above, font=\scriptsize] {$\mu = 6$};
  \node[font=\scriptsize, anchor=south] at (axis cs:1.9,0.2) {$\mu - 2\sigma$};
  \node[font=\scriptsize, anchor=south] at (axis cs:10.1,0.2) {$\mu + 2\sigma$};
  \node[font=\small, anchor=west, text width=5.2cm, align=left] at (axis cs:11.5,0.13)
    {$\mathcal{B}(20;\ 0{,}3)$: $\sigma = \sqrt{4{,}2} \approx 2{,}05$.\\ Nella fascia gialla c'è il 97,5\% della
     probabilità; Chebyshev ne garantisce almeno il 75\%.};
\end{axis}
\end{tikzpicture}
\caption{La fascia $\mu \pm 2\sigma$ per una binomiale. Come Markov, anche Chebyshev è prudente: vale per tutte le pmf,
quindi per una pmf particolare la probabilità vera è di solito molto più alta del limite garantito.}
\end{figure}

Riscrivendo il teorema in termini relativi, il parametro fondamentale risulta il rapporto $\mu_X/\sigma_X$. Valori
elevati di questo rapporto indicano una pmf molto concentrata intorno alla media rispetto alla sua posizione (una
variabile «poco aleatoria»), mentre valori bassi indicano un'elevata aleatorietà. (Una funzione è molto concentrata se
è sostanzialmente diversa da zero in un intervallo piccolo rispetto al suo valore centrale.)

\begin{esempio}{Poisson grandi e piccole}
Per una Poisson $\mu_X/\sigma_X = \lambda/\sqrt{\lambda} = \sqrt{\lambda}$. Con $\lambda = 1$ il rapporto vale 1: la
variabile vale spesso 0, 1, 2, 3, cioè fluttua quanto è grande. Con $\lambda = 100$ il rapporto vale 10: la variabile
vale circa $100 \pm 10$, fluttuazioni del 10\%. Con $\lambda = 10\,000$, circa $10\,000 \pm 100$: l'1\%. Più eventi
rari si contano, più il loro numero diventa prevedibile in proporzione.
\end{esempio}

\gbox{In sintesi}{mspiA!80!black}{
\begin{itemize}
  \item Valore quadratico medio $E[X^2]$, valore efficace $\sqrt{E[X^2]}$, varianza
        $\sigma_X^2 = E[(X - \mu_X)^2] = E[X^2] - \mu_X^2$, deviazione standard $\sigma_X$.
  \item Bernoulli: $p(1-p)$. Binomiale: $N p (1-p)$. Uniforme su $\{0, \ldots, M-1\}$: $(M^2-1)/12$. Poisson:
        $\lambda$ (uguale alla media). Si calcolano con il trucco della derivata o con $E[X(X-1)]$.
  \item La \textbf{derivata} è la pendenza della tangente: $(x^n)' = n x^{n-1}$, $(fg)' = f'g + fg'$, $(e^x)' = e^x$; da
        qui la serie $e^x = \sum x^k/k!$.
  \item $E[aX + b] = aE[X] + b$; $\sigma_{aX+b}^2 = a^2 \sigma_X^2$; $E[X] \geq 0$ se $X \geq 0$; $\sigma^2 \geq 0$.
  \item \textbf{Markov}: $P(Z \geq \delta) \leq E[Z]/\delta$ per $Z \geq 0$. \textbf{Chebyshev}:
        $P(\abs{X - \mu_X} \geq k\sigma_X) \leq 1/k^2$, cioè $P(\mu_X - k\sigma_X \leq X \leq \mu_X + k\sigma_X) \geq
        1 - 1/k^2$: la deviazione standard è il righello della dispersione, e $\mu/\sigma$ misura la concentrazione.
\end{itemize}
}
